% Problemas del capítulo 4. Generado por tools/problemas/gen.py; no editar a mano.
\section{Problemas}\label{sec:problemas-04}

Graduados: \({\star}\) directo, \({\star}{\star}\) para pensar, \({\star}{\star}{\star}\) desafío. Las respuestas están en el apéndice~\ref{ap:respuestas}; conviene intentarlos antes de mirarlas.

\begin{problema}[\({\star}\)]\label{pb:04-1}
Encontrá la afinidad \(x\mapsto Ax+b\) que lleva \((0,0)\), \((1,0)\) y \((0,1)\) a \((2,1)\), \((5,2)\) y \((3,4)\). ¿A dónde va \((1,1)\)? ¿Por cuánto multiplica las áreas?
\end{problema}

\begin{problema}[\({\star}{\star}\)]\label{pb:04-2}
Con la afinidad del problema anterior, comprobá que el punto medio de \((0,0)\) y \((2,2)\) va al punto medio de las imágenes. ¿Pasa lo mismo con los ángulos?
\end{problema}

\begin{problema}[\({\star}{\star}\)]\label{pb:04-3}
Un plano escaneado se encogió un \(0{,}4\)\,\% en sentido horizontal y nada en vertical. Escribí la afinidad que lo corrige. Un lote mide \(30\times20\)~m sobre el plano escaneado (\(30\) en horizontal): ¿cuál es su superficie real?
\end{problema}

\begin{problema}[\({\star}{\star}\)]\label{pb:04-4}
En coordenadas homogéneas, escribí las matrices de «trasladar \((2,3)\)» y de «escalar \(\times2\) desde el origen». Aplicá a \((1,1)\) las dos en un orden y en el otro.
\end{problema}

\begin{problema}[\({\star}{\star}{\star}\)]\label{pb:04-5}
La afinidad \(x\mapsto Ax+b\), con \(A=\begin{pmatrix}0&-0{,}5\\0{,}5&0\end{pmatrix}\) y \(b=(3,1)\), gira \(90^\circ\) y achica a la mitad. Encontrá su punto fijo, resolviendo \((I-A)p=b\).
\end{problema}
